\chapter*{Resumo}
A experiência musical envolve dimensões culturais e sociais que nem sempre são igualmente acessíveis a pessoas surdas ou com deficiência auditiva, especialmente quando as interfaces dependem predominantemente da audição. Nesse contexto, este trabalho apresenta o desenvolvimento e a validação técnica de uma interface musical digital multissensorial voltada à acessibilidade musical. A arquitetura utiliza o controlador MIDI Akai MPD218 como interface de entrada e emprega o parâmetro \textit{velocity} como representação digital da intensidade do gesto realizado sobre os \textit{pads}. Os eventos MIDI são processados em Python, convertidos em comandos PWM e transmitidos por UDP sobre rede Wi-Fi a uma unidade embarcada baseada no ESP32. Na configuração final, a unidade vestível integra um motor vibrotátil Vibracall 1027, acionado por estágio com transistor MOSFET, e um LED de alto brilho, além de alimentação própria composta por duas células de íons de lítio em configuração 2S, módulo BMS e conversor DC--DC. A validação envolveu simulações do estágio de acionamento, testes de comunicação, ensaios da lógica de controle e caracterizações experimentais da entrada e do atuador. No teste UDP, os 100 pacotes transmitidos foram registrados pelo ESP32, sem perdas na sequência avaliada. No ensaio com corpo de 39~g, foram realizadas 60 quedas entre 0,5~cm e 5~cm; a associação entre altura e \textit{velocity} apresentou coeficiente de Spearman de 0,944 ($p<0{,}001$). Na caracterização acústica do motor, seis níveis de PWM entre 80 e 255 foram avaliados em três repetições, com frequências médias entre 91,5~Hz e 227,9~Hz e coeficiente de Spearman de 0,988 ($p<0{,}001$). Com alimentação embarcada, foram medidos valores entre aproximadamente 3,26~V e 3,29~V no barramento do ESP32 com Wi-Fi ativo. A operação integrada confirmou a cadeia funcional entre interação MIDI, processamento computacional, comunicação sem fio e respostas vibrotátil e visual. Os resultados demonstram a viabilidade técnica da arquitetura como plataforma experimental, sem estabelecer, nesta etapa, eficácia perceptiva para pessoas surdas ou com deficiência auditiva.

\noindent\textbf{Palavras-chave:} acessibilidade musical; MIDI; \textit{feedback} vibrotátil; interação multissensorial; ESP32.

\clearpage
